%============================================================
% CH1 · P5 · PROBLEM 1 — HOW MANY NUMBERS ARE SMALLER THAN THE CURRENT NUMBER (LeetCode 1365)
%============================================================
\subsection{Problem 1.5.1: How Many Numbers Are Smaller Than the Current Number}
\problemlabel{1.5.1}{ssec:ch1-smaller-than-current}

\problemcard{1.5.1}{How Many Numbers Are Smaller Than the Current Number}{LeetCode 1365}{https://leetcode.com/problems/how-many-numbers-are-smaller-than-the-current-number/}{\exeasy\ \lc{Easy}}{Counting array over the values $0..100$, then prefix counts: \texttt{less[v]} = number of elements $< v$}

\subsubsection*{1. Problem statement}

Given an array \texttt{nums}, return the array \texttt{ans} where $\texttt{ans}[i]$ is the number of indices $j$ with $\texttt{nums}[j] < \texttt{nums}[i]$ (strictly smaller; $j \ne i$ is automatic).

\textbf{Constraints:} $2 \le n \le 500$, \; $0 \le \texttt{nums[i]} \le 100$.

\begin{table}[H]
\centering\small
\begin{tabular}{@{}lll@{}}
\toprule
\textbf{nums} & \textbf{Output} & \textbf{Why} \\
\midrule
$[8,1,2,2,3]$ & $[4,0,1,1,3]$ & four values are below 8; only the 1 is below each 2 \\
$[6,5,4,8]$ & $[2,1,0,3]$ & distinct values: the answer is the rank in sorted order \\
$[7,7,7,7]$ & $[0,0,0,0]$ & equal values are not smaller \\
\bottomrule
\end{tabular}
\end{table}

\begin{diagrambox}
\begin{tikzpicture}[x=1cm, y=1cm]
    \dpindex{ix}{1.75}{index}{0}{4}
    \dparray{a}{1.1}{nums}{8/aread, 1/acell, 2/acell, 2/acell, 3/acell}
    \dparray{b}{-0.3}{ans}{4/awrite, 0/akeep, 1/akeep, 1/akeep, 3/akeep}
    \foreach \j in {1,2,3,4}{ \draw[faded] (a\j.north) to[bend right=30] (a0.north); }
    \node[note, anchor=west, text width=4.8cm, align=left] at (5.2,0.4) {$\texttt{ans}[0] = 4$: all four other values $1, 2, 2, 3$ are $< 8$. $\texttt{ans}[2] = \texttt{ans}[3] = 1$: the two 2s do not count each other.};
\end{tikzpicture}
\figcap{Sample 1. Equal values get equal answers.}
\end{diagrambox}

\subsubsection*{2. The problem model}
\begin{modelbox}[Problem Model --- Smaller Than the Current Number]
\renewcommand{\arraystretch}{1.35}
\begin{tabularx}{\linewidth}{@{}>{\bfseries}l X@{}}
Input & an array $a[0..n-1]$ with values in $[0, 100]$ \\
Computed & an array: $\texttt{ans}[i] = |\{\, j : a_j < a_i \,\}|$, the \emph{rank} of $a_i$ (0-based, ties share the smallest rank) \\
Objects & pairs $(i, j)$ --- but the answer depends on $a_i$ only through its value, so it is a function $\texttt{less}(v)$ of the value \\
Size & $n \le 500$: even $\bigO{n^2} = 2.5\cdot10^5$ passes; the value range $V = 101$ invites $\bigO{n + V}$ \\
\end{tabularx}
\end{modelbox}

\begin{diagrambox}
\begin{tikzpicture}[x=1cm, y=1cm]
    \draw[->, thick, darkgray] (-0.2,0) -- (9.6,0) node[right, font=\scriptsize] {value};
    \foreach \v in {0,1,2,3,4,5,6,7,8,9}{ \draw (\v,0.08) -- (\v,-0.08) node[below, font=\scriptsize] {\v}; }
    \foreach \v/\h in {1/1, 2/2, 3/1, 8/1}{
        \foreach \k in {1,...,\h}{ \node[circle, fill=primaryblue, inner sep=2.2pt] at (\v, 0.12 + 0.3*\k) {}; }
    }
    \draw[decorate, decoration={brace, amplitude=5pt}, thick, memoorange] (0.8,1.0) -- (3.2,1.0) node[midway, above=5pt, font=\scriptsize\bfseries, text=memoorange] {4 values $< 8$};
    \draw[dep] (8,1.3) -- (8,0.55);
    \node[font=\scriptsize\bfseries, text=accentorange] at (8,1.5) {$\texttt{less}(8) = 4$};
\end{tikzpicture}
\figcap{The array as dots on the value axis (Sample 1). The answer for value $v$ is the number of dots strictly left of $v$ --- a prefix sum over the value axis.}
\end{diagrambox}

\subsubsection*{3. Key observations}

\begin{itemize}[leftmargin=*]
    \item \textbf{The answer depends only on the value.} Equal values get equal answers, so we need $\texttt{less}(v)$ for at most 101 values, not $n$ separate counts.
    \item \textbf{$\texttt{less}(v)$ is a prefix sum of counts:} $\texttt{less}(v) = \sum_{u < v} \texttt{cnt}[u]$. One pass to count, one pass over 101 values to accumulate.
    \item \textbf{Tempting wrong idea --- ``the position in the sorted array''.} Sort the indices by value and answer the position of each index. Without duplicates (Sample 2) this is right, but on Sample 1 the sorted order is $1, 2, 2, 3, 8$ and the \emph{second} 2 sits at position 2, while its answer is 1. Ties must share the position of the \emph{first} copy.
    \item \textbf{Another wrong idea --- ``(count of values $\le v$) $- 1$''.} It subtracts the element itself but not its equal twins: for the 2s of Sample 1 it gives $3 - 1 = 2$ instead of 1.
\end{itemize}

\begin{diagrambox}
\begin{tikzpicture}[x=1cm, y=1cm]
    \dparray{s}{1.1}{sorted}{1/acell, 2/akeep, 2/aout, 3/acell, 8/acell}
    \aptr{s2}{second 2 at position 2: wrong}{warningred}
    \aptrup{s1}{first copy at position 1 \yes}{successgreen}
    \node[note, anchor=west, text width=5.0cm, align=left] at (5.2,1.1) {the position of an element in sorted order counts the equal elements placed before it; the answer must not};
\end{tikzpicture}
\figcap{Counter-example to ``answer = position after sorting'': only the first copy of each value is at the right position.}
\end{diagrambox}

\subsubsection*{4. Approaches}

\begin{bruteforcebox}
\textbf{Idea.} For every $i$, scan the whole array and count the smaller values.
\begin{lstlisting}[style=cpp]
#include <bits/stdc++.h>
using namespace std;

class Solution {
public:
    vector<int> smallerNumbersThanCurrent(vector<int>& nums) {
        int n = nums.size();
        vector<int> ans(n, 0);
        for (int i = 0; i < n; i++)
            for (int j = 0; j < n; j++)
                if (nums[j] < nums[i]) ans[i]++;   // j == i never counts
        return ans;
    }
};

int main() {                                       // input: n, then the n values
    int n;
    if (!(cin >> n)) return 0;
    vector<int> a(n);
    for (auto& x : a) cin >> x;
    vector<int> r = Solution().smallerNumbersThanCurrent(a);
    for (int i = 0; i < n; i++) cout << r[i] << (i + 1 < n ? ' ' : '\n');
}
\end{lstlisting}
\begin{complexitybox}
\textbf{Time:} $\bigO{n^2}$ --- $2.5\cdot10^5$ comparisons for $n = 500$: \ac{} here, but hopeless for $n = 10^5$ ($10^{10}$). \textbf{Space:} $\bigO{1}$ besides the answer.
\end{complexitybox}
\end{bruteforcebox}

\begin{algorithmbox}
\textbf{Idea.} Count every value (Tool~1), turn the counts into prefix counts (Tool~3), and answer $\texttt{less}[\texttt{nums}[i]]$ for each $i$.
\begin{lstlisting}[style=cpp]
#include <bits/stdc++.h>
using namespace std;

class Solution {
public:
    vector<int> smallerNumbersThanCurrent(vector<int>& nums) {
        const int MAXV = 100;
        vector<int> cnt(MAXV + 2, 0);             // cnt[v] = how many times v occurs
        for (int x : nums) cnt[x]++;
        vector<int> less(MAXV + 2, 0);            // less[v] = how many values are < v
        for (int v = 1; v <= MAXV + 1; v++) less[v] = less[v - 1] + cnt[v - 1];
        vector<int> ans;
        for (int x : nums) ans.push_back(less[x]);
        return ans;
    }
};

int main() {                                      // input: n, then the n values (0..100)
    int n;
    if (!(cin >> n)) return 0;
    vector<int> a(n);
    for (auto& x : a) cin >> x;
    vector<int> r = Solution().smallerNumbersThanCurrent(a);
    for (int i = 0; i < n; i++) cout << r[i] << (i + 1 < n ? ' ' : '\n');
}
\end{lstlisting}
\begin{complexitybox}
\textbf{Time:} $\bigO{n + V}$ with $V = 101$ --- three short loops. \textbf{Space:} $\bigO{V}$ for \texttt{cnt} and \texttt{less}. \textbf{Verdict:} \ac{}.
\end{complexitybox}
\end{algorithmbox}

\begin{invariantbox}[Why \texttt{less[v]} is the answer]
After the counting loop, $\texttt{cnt}[u]$ is the number of elements equal to $u$ (invariant of Tool~1). The accumulation keeps $\texttt{less}[v] = \sum_{u=0}^{v-1} \texttt{cnt}[u]$: true for $v = 0$ (empty sum), and $\texttt{less}[v] = \texttt{less}[v-1] + \texttt{cnt}[v-1]$ adds exactly the elements equal to $v - 1$. Since values are integers, ``$a_j < v$'' means ``$a_j = u$ for some $u \le v-1$'', and these events are disjoint for different $u$. Hence $\texttt{less}[a_i] = |\{ j : a_j < a_i \}| = \texttt{ans}[i]$. \hfill$\blacksquare$
\end{invariantbox}

\begin{alternativebox}
\textbf{Sort a copy and binary-search the first copy.} In sorted order the first copy of $v$ sits at index $\texttt{less}(v)$. So sort a copy of the array and answer \texttt{lower\_bound(s, nums[i]) - s.begin()} --- the index of the first element $\ge \texttt{nums}[i]$, which is the number of elements $< \texttt{nums}[i]$. This does not depend on the value range at all.
\begin{lstlisting}[style=cpp]
#include <bits/stdc++.h>
using namespace std;

class Solution {
public:
    vector<int> smallerNumbersThanCurrent(vector<int>& nums) {
        vector<int> s(nums);
        sort(s.begin(), s.end());
        vector<int> ans;
        for (int x : nums)                        // index of the first copy of x in sorted order
            ans.push_back(lower_bound(s.begin(), s.end(), x) - s.begin());
        return ans;
    }
};

int main() {                                      // input: n, then the n values
    int n;
    if (!(cin >> n)) return 0;
    vector<int> a(n);
    for (auto& x : a) cin >> x;
    vector<int> r = Solution().smallerNumbersThanCurrent(a);
    for (int i = 0; i < n; i++) cout << r[i] << (i + 1 < n ? ' ' : '\n');
}
\end{lstlisting}
$\bigO{n \log n}$ time, $\bigO{n}$ memory. It is the method to use when values go up to $10^9$ (variant V3); binary search itself is the topic of Chapter~\ref{ch:binary_search}.

\begin{diagrambox}
\begin{tikzpicture}[x=1cm, y=1cm]
    \dpindex{qx}{0.65}{index}{0}{4}
    \dparray{q}{0}{sorted}{1/acell, 2/aread, 2/awin, 3/acell, 8/acell}
    \aptr{q1}{lower\_bound(2) = 1}{accentorange}
    \node[note, anchor=west, text width=4.4cm, align=left] at (5.2,0) {the first element $\ge 2$ is at index 1, so exactly one element is $< 2$};
\end{tikzpicture}
\figcap{\texttt{lower\_bound} lands on the first copy, which is exactly what ties need.}
\end{diagrambox}
\end{alternativebox}

\begin{tricksbox}
\begin{itemize}[leftmargin=*]
    \item \textbf{One array is enough:} turn \texttt{cnt} into \texttt{less} in place by a running sum that \emph{excludes} the current cell: \texttt{for v: \{ int c = cnt[v]; cnt[v] = run; run += c; \}}.
    \item \textbf{Size $\text{MAX} + 2$:} \texttt{less} needs index $\text{MAX} + 1$ only if you also want ``how many are $\le \text{MAX}$''; allocating one extra cell removes the boundary check.
    \item \textbf{Strictly smaller vs.\ not larger:} ``$\le$ excluding itself'' is $\texttt{less}[v+1] - 1$.
\end{itemize}
\begin{diagrambox}
\begin{tikzpicture}[x=1cm, y=1cm]
    \dpindex{w}{1.35}{$v$}{0}{4}
    \dparray{p}{0.7}{\texttt{cnt}}{0/compcell, 1/compcell, 2/compcell, 1/compcell, 0/compcell}
    \dparray{q}{-0.6}{in place}{0/basecell, 0/basecell, 1/basecell, 3/basecell, 4/basecell}
    \foreach \k in {0,...,4}{ \draw[faded] (p\k.south) -- (q\k.north); }
    \node[note, anchor=west, text width=5.2cm, align=left] at (4.8,0.05) {\texttt{run} before each cell: $0, 0, 1, 3, 4$; each cell receives \texttt{run} and then adds its old count to it (Sample 1, values $0..4$)};
\end{tikzpicture}
\figcap{Turning \texttt{cnt} into \texttt{less} in place with a running sum that excludes the current cell.}
\end{diagrambox}
\end{tricksbox}

\subsubsection*{5. Trace}

\begin{tracebox}[{Sample 1: $[8,1,2,2,3]$}]
\centering
\begin{tikzpicture}[x=1cm, y=1cm]
    \dparray{a}{2.6}{nums}{8/acell, 1/acell, 2/acell, 2/acell, 3/acell}
    \dpindex{v}{1.6}{value $v$}{0}{9}
    \dparray{c}{0.95}{\texttt{cnt}}{0/compcell, 1/compcell, 2/compcell, 1/compcell, 0/compcell, 0/compcell, 0/compcell, 0/compcell, 1/compcell, 0/compcell}
    \dparray{l}{-0.35}{\texttt{less}}{0/basecell, 0/basecell, 1/basecell, 3/basecell, 4/basecell, 4/basecell, 4/basecell, 4/basecell, 4/anscell, 5/basecell}
    \foreach \k in {0,...,8}{ \pgfmathtruncatemacro{\kk}{\k+1} \draw[faded] (c\k.south east) -- (l\kk.north west); }
    \dparray{r}{-1.75}{\texttt{ans}}{4/akeep, 0/akeep, 1/akeep, 1/akeep, 3/akeep}
\end{tikzpicture}
\figcap{Counting, then accumulating: each \texttt{less} cell is the cell to its left plus the count diagonally above-left (thin lines). The answer reads \texttt{less} at each element's value; the orange cell is $\texttt{less}[8] = 4$.}

\medskip
\small
\begin{tabular}{@{}lcccccccccc@{}}
\toprule
$v$ & 0 & 1 & 2 & 3 & 4 & 5 & 6 & 7 & 8 & 9 \\
\midrule
\texttt{cnt[v]} & 0 & 1 & 2 & 1 & 0 & 0 & 0 & 0 & 1 & 0 \\
\texttt{less[v]} & 0 & 0 & 1 & 3 & 4 & 4 & 4 & 4 & 4 & 5 \\
\bottomrule
\end{tabular}

\medskip
\begin{tabular}{@{}lccccc@{}}
\toprule
$i$ & 0 & 1 & 2 & 3 & 4 \\
\midrule
\texttt{nums[i]} & 8 & 1 & 2 & 2 & 3 \\
\texttt{ans[i]} $= \texttt{less}[\texttt{nums}[i]]$ & 4 & 0 & 1 & 1 & 3 \\
\bottomrule
\end{tabular}
\end{tracebox}

\subsubsection*{6. Test suite}

\begin{testbox}
\centering\small
\begin{tabular}{@{}lll@{}}
\toprule
\textbf{nums} & \textbf{Why it is interesting} & \textbf{Expected} \\
\midrule
$[8,1,2,2,3]$ & official Sample 1: duplicates & $[4,0,1,1,3]$ \\
$[6,5,4,8]$ & official Sample 2: distinct values & $[2,1,0,3]$ \\
$[7,7,7,7]$ & official Sample 3: all equal & $[0,0,0,0]$ \\
$[0,100]$ & both ends of the value range & $[0,1]$ \\
$[100,100,0]$ & largest value repeated & $[1,1,0]$ \\
$[5,4,3,2,1]$ & decreasing & $[4,3,2,1,0]$ \\
\bottomrule
\end{tabular}

\medskip
The three programs agree on the official samples and on 500 random arrays ($2 \le n \le 10$, values drawn from small ranges and from $\{0, 100\}$ to hit the borders).
\end{testbox}

\subsubsection*{7. Common mistakes}

\begin{warningbox}[Mistake 1 --- the counting array one cell too short]
Values go up to 100 inclusive: \texttt{vector<int> cnt(100)} has no cell for 100. The test $[0,100]$ catches it.
\end{warningbox}

\begin{warningbox}[Mistake 2 --- including the current cell in the prefix]
\begin{lstlisting}[style=snippet]
less[v] = less[v - 1] + cnt[v];       // WRONG: counts values <= v (the 2s get 3)
less[v] = less[v - 1] + cnt[v - 1];   // RIGHT: values < v
\end{lstlisting}
\end{warningbox}

\begin{warningbox}[Mistake 3 --- using sorted positions with ties]
Sorting indices by value and taking each index's position gives $[4,0,1,2,3]$ on Sample 1. Ties must take the position of the first copy (\texttt{lower\_bound}), or use the counts.
\end{warningbox}

\subsubsection*{8. Takeaways}

\begin{itemize}[leftmargin=*]
    \item If the answer depends only on the value, compute it once per value, not once per element.
    \item ``How many are smaller than $v$'' is a prefix sum over the counting array; ``how many are larger'' is $n$ minus a prefix sum.
    \item Prefix counts are the starts of the value blocks in sorted order --- the heart of counting sort.
    \item With large values the same ranks come from sorting plus \texttt{lower\_bound}.
\end{itemize}

\begin{variantsbox}
\variant{Count Elements With Strictly Smaller and Greater Elements}{\exeasy\ \lc{Easy}\ \varlink{Count Elements With Strictly Smaller and Greater Elements}{LeetCode 2148}{https://leetcode.com/problems/count-elements-with-strictly-smaller-and-greater-elements/}}
  {count the elements that have at least one strictly smaller and one strictly larger element.}
  {model: $\texttt{less}(a_i) > 0$ and $n - \texttt{less}(a_i + 1) > 0$ $\iff \min < a_i < \max$.}
  {algorithm: one pass for $\min$ and $\max$, one to count; values up to $10^5$ in absolute value. Samples $\to 2, 2$.}
\variant{Height Checker}{\exeasy\ \lc{Easy}\ \varlink{Height Checker}{LeetCode 1051}{https://leetcode.com/problems/height-checker/}}
  {count the positions where the array differs from its sorted version (heights $1..100$).}
  {model: the sorted array is fully described by \texttt{cnt}: counting sort.}
  {algorithm: walk the array and a pointer $v$ over the values with $\texttt{cnt}[v] > 0$ at the same time; $\bigO{n + 100}$. Samples $\to 3, 5, 0$.}
\variant{Large values}{\varnone}
  {the same question with $|a_i| \le 10^9$ and $n \le 10^5$.}
  {model: the value range is too large for an array; ranks are still positions of first copies in sorted order.}
  {algorithm: the alternative (sort + \texttt{lower\_bound}), $\bigO{n \log n}$; or compress values first.}
\variant{How many are larger}{\varnone}
  {$\texttt{ans}[i] = |\{ j : a_j > a_i \}|$.}
  {model: complement of ``$\le$'': $n - \texttt{less}(a_i + 1)$.}
  {algorithm: the same prefix array, read one cell further; or accumulate from the top value downwards.}
\end{variantsbox}
